% ═══════════════════════════════════════════════════════════════
% 民事起诉状 · 共享调解意愿 + 具状人签名
% 通过 % ═══════════════════════════════════════════════════════════════
% 民事起诉状 · 共享调解意愿 + 具状人签名
% 通过 % ═══════════════════════════════════════════════════════════════
% 民事起诉状 · 共享调解意愿 + 具状人签名
% 通过 % ═══════════════════════════════════════════════════════════════
% 民事起诉状 · 共享调解意愿 + 具状人签名
% 通过 \input{../common/mediation} 引入
% ═══════════════════════════════════════════════════════════════

% 对纠纷解决方式的意愿副标题
\subtitlesection{对纠纷解决方式的意愿}

% 是否了解调解作为非诉讼纠纷解决方式
\pairsectionleft{是否了解调解作为非诉\\讼纠纷解决方式，能\\及时、高效、低成本、\\不伤和气地解决纠纷}{
了解$\square$\hspace{1em}不了解$\square$
}

% 是否了解先行调解解决纠纷的好处
\pairsectionleft{是否了解先行调解解\\决纠纷的好处}{
1. 立案后选择先行调解的，可以很快启动调解程序。如不同意调解，法院\\
将依程序开庭审理案件，但可能需要经过较长一段时间的排期等待，且审\\
理、执行周期相对较长。\\
了解$\square$\hspace{1em}不了解$\square$\\
2. 选择先行调解，调解成功且自动履行的免交诉讼费用，申请司法确认的\\
不交纳诉讼费用，要求出具调解书的减半交纳诉讼费用。\\
了解$\square$\hspace{1em}不了解$\square$\\
3. 首次调解不成功，但仍有继续调解意愿的，可以选择更换调解组织和调\\
解员再进行调解。调解无法达成一致意见的，法院将依程序排期开庭。\\
了解$\square$\hspace{1em}不了解$\square$\\
4. 依照法律规定，调解具有保密性要求，调解过程不公开，调解协议未经\\
当事人同意不得公开。\\
了解$\square$\hspace{1em}不了解$\square$\\
5. 调解达成的协议具有法律效力，可以依照法律规定申请司法确认，具有\\
强制执行效力。\\
了解$\square$\hspace{1em}不了解$\square$
}

% 是否考虑先行调解
\pairsectionleft{是否考虑先行调解}{
是$\square$\\
否$\square$\\
暂不确定，想要了解更多内容$\square$
}

% 具状人签字与日期
\vspace{10bp}
\noindent\hspace*{248bp}{\fzdbs\fontsize{15bp}{22bp}\selectfont 具状人（签字、盖章）：}\par
\noindent\hspace*{248bp}{\fzdbs\fontsize{15bp}{22bp}\selectfont 日期：}
 引入
% ═══════════════════════════════════════════════════════════════

% 对纠纷解决方式的意愿副标题
\subtitlesection{对纠纷解决方式的意愿}

% 是否了解调解作为非诉讼纠纷解决方式
\pairsectionleft{是否了解调解作为非诉\\讼纠纷解决方式，能\\及时、高效、低成本、\\不伤和气地解决纠纷}{
了解$\square$\hspace{1em}不了解$\square$
}

% 是否了解先行调解解决纠纷的好处
\pairsectionleft{是否了解先行调解解\\决纠纷的好处}{
1. 立案后选择先行调解的，可以很快启动调解程序。如不同意调解，法院\\
将依程序开庭审理案件，但可能需要经过较长一段时间的排期等待，且审\\
理、执行周期相对较长。\\
了解$\square$\hspace{1em}不了解$\square$\\
2. 选择先行调解，调解成功且自动履行的免交诉讼费用，申请司法确认的\\
不交纳诉讼费用，要求出具调解书的减半交纳诉讼费用。\\
了解$\square$\hspace{1em}不了解$\square$\\
3. 首次调解不成功，但仍有继续调解意愿的，可以选择更换调解组织和调\\
解员再进行调解。调解无法达成一致意见的，法院将依程序排期开庭。\\
了解$\square$\hspace{1em}不了解$\square$\\
4. 依照法律规定，调解具有保密性要求，调解过程不公开，调解协议未经\\
当事人同意不得公开。\\
了解$\square$\hspace{1em}不了解$\square$\\
5. 调解达成的协议具有法律效力，可以依照法律规定申请司法确认，具有\\
强制执行效力。\\
了解$\square$\hspace{1em}不了解$\square$
}

% 是否考虑先行调解
\pairsectionleft{是否考虑先行调解}{
是$\square$\\
否$\square$\\
暂不确定，想要了解更多内容$\square$
}

% 具状人签字与日期
\vspace{10bp}
\noindent\hspace*{248bp}{\fzdbs\fontsize{15bp}{22bp}\selectfont 具状人（签字、盖章）：}\par
\noindent\hspace*{248bp}{\fzdbs\fontsize{15bp}{22bp}\selectfont 日期：}
 引入
% ═══════════════════════════════════════════════════════════════

% 对纠纷解决方式的意愿副标题
\subtitlesection{对纠纷解决方式的意愿}

% 是否了解调解作为非诉讼纠纷解决方式
\pairsectionleft{是否了解调解作为非诉\\讼纠纷解决方式，能\\及时、高效、低成本、\\不伤和气地解决纠纷}{
了解$\square$\hspace{1em}不了解$\square$
}

% 是否了解先行调解解决纠纷的好处
\pairsectionleft{是否了解先行调解解\\决纠纷的好处}{
1. 立案后选择先行调解的，可以很快启动调解程序。如不同意调解，法院\\
将依程序开庭审理案件，但可能需要经过较长一段时间的排期等待，且审\\
理、执行周期相对较长。\\
了解$\square$\hspace{1em}不了解$\square$\\
2. 选择先行调解，调解成功且自动履行的免交诉讼费用，申请司法确认的\\
不交纳诉讼费用，要求出具调解书的减半交纳诉讼费用。\\
了解$\square$\hspace{1em}不了解$\square$\\
3. 首次调解不成功，但仍有继续调解意愿的，可以选择更换调解组织和调\\
解员再进行调解。调解无法达成一致意见的，法院将依程序排期开庭。\\
了解$\square$\hspace{1em}不了解$\square$\\
4. 依照法律规定，调解具有保密性要求，调解过程不公开，调解协议未经\\
当事人同意不得公开。\\
了解$\square$\hspace{1em}不了解$\square$\\
5. 调解达成的协议具有法律效力，可以依照法律规定申请司法确认，具有\\
强制执行效力。\\
了解$\square$\hspace{1em}不了解$\square$
}

% 是否考虑先行调解
\pairsectionleft{是否考虑先行调解}{
是$\square$\\
否$\square$\\
暂不确定，想要了解更多内容$\square$
}

% 具状人签字与日期
\vspace{10bp}
\noindent\hspace*{248bp}{\fzdbs\fontsize{15bp}{22bp}\selectfont 具状人（签字、盖章）：}\par
\noindent\hspace*{248bp}{\fzdbs\fontsize{15bp}{22bp}\selectfont 日期：}
 引入
% ═══════════════════════════════════════════════════════════════

% 对纠纷解决方式的意愿副标题
\subtitlesection{对纠纷解决方式的意愿}

% 是否了解调解作为非诉讼纠纷解决方式
\pairsectionleft{是否了解调解作为非诉\\讼纠纷解决方式，能\\及时、高效、低成本、\\不伤和气地解决纠纷}{
了解$\square$\hspace{1em}不了解$\square$
}

% 是否了解先行调解解决纠纷的好处
\pairsectionleft{是否了解先行调解解\\决纠纷的好处}{
1. 立案后选择先行调解的，可以很快启动调解程序。如不同意调解，法院\\
将依程序开庭审理案件，但可能需要经过较长一段时间的排期等待，且审\\
理、执行周期相对较长。\\
了解$\square$\hspace{1em}不了解$\square$\\
2. 选择先行调解，调解成功且自动履行的免交诉讼费用，申请司法确认的\\
不交纳诉讼费用，要求出具调解书的减半交纳诉讼费用。\\
了解$\square$\hspace{1em}不了解$\square$\\
3. 首次调解不成功，但仍有继续调解意愿的，可以选择更换调解组织和调\\
解员再进行调解。调解无法达成一致意见的，法院将依程序排期开庭。\\
了解$\square$\hspace{1em}不了解$\square$\\
4. 依照法律规定，调解具有保密性要求，调解过程不公开，调解协议未经\\
当事人同意不得公开。\\
了解$\square$\hspace{1em}不了解$\square$\\
5. 调解达成的协议具有法律效力，可以依照法律规定申请司法确认，具有\\
强制执行效力。\\
了解$\square$\hspace{1em}不了解$\square$
}

% 是否考虑先行调解
\pairsectionleft{是否考虑先行调解}{
是$\square$\\
否$\square$\\
暂不确定，想要了解更多内容$\square$
}

% 具状人签字与日期
\vspace{10bp}
\noindent\hspace*{248bp}{\fzdbs\fontsize{15bp}{22bp}\selectfont 具状人（签字、盖章）：}\par
\noindent\hspace*{248bp}{\fzdbs\fontsize{15bp}{22bp}\selectfont 日期：}
